% Part:sets-functions-relations
% Chapter: size-of-sets
% Section: zig-zag

\documentclass[../../../include/open-logic-section]{subfiles}

\begin{document}

\olfileid{sfr}{siz}{zigzag}

\olsection{De zigzagmethode van Cantor}

\begin{explain}
We hebben al enkele ``eenvoudige'' opsommingen bekeken. Nu volgt een
iets moeilijker geval. Beschouw de verzameling paren van natuurlijke
getallen\oliflabeldef{sfr:set:pai:sec}{, die we in
\olref[set][pai]{sec} als volgt hebben gedefinieerd:}{, gedefinieerd door:}
\[
\Nat \times \Nat = \Setabs{\tuple{n,m}}{n,m \in \Nat}
\]
We kunnen deze geordende paren als volgt in een \emph{tabel} rangschikken:
\[
\begin{array}{ c | c | c | c | c | c}
& \mathbf 0 & \mathbf 1 & \mathbf 2 & \mathbf 3 & \dots \\
\hline
\mathbf 0 & \tuple{0,0} & \tuple{0,1} & \tuple{0,2} & \tuple{0,3} & \dots \\
\hline
\mathbf 1 & \tuple{1,0} & \tuple{1,1} & \tuple{1,2} & \tuple{1,3} & \dots \\
\hline
\mathbf 2 & \tuple{2,0} & \tuple{2,1} & \tuple{2,2} & \tuple{2,3} & \dots \\
\hline
\mathbf 3 & \tuple{3,0} & \tuple{3,1} & \tuple{3,2} & \tuple{3,3} & \dots \\
\hline
\vdots & \vdots & \vdots & \vdots & \vdots & \ddots\\
\end{array}
\]
Elk geordend paar uit $\Nat \times \Nat$ komt duidelijk precies eenmaal in de
tabel voor. Het paar $\tuple{n,m}$ staat in het bijzonder in rij $n$ en kolom
$m$. Maar hoe rangschikken we de elementen van zo'n tabel in een
``eendimensionale'' lijst? Het patroon in de volgende tabel geeft één
mogelijkheid (er zijn uiteraard vele andere):
\[
\begin{array}{ c | c | c | c | c | c | c}
& \mathbf 0 & \mathbf 1 & \mathbf 2 & \mathbf 3 & \mathbf 4 &\dots \\
\hline
\mathbf 0 & 0  & 1& 3 & 6& 10 &\ldots \\
\hline
\mathbf 1 &2 & 4& 7 & 11 & \dots &\ldots \\
\hline
\mathbf 2 & 5 & 8 & 12 & \ldots & \dots&\ldots \\
\hline
\mathbf 3 & 9 & 13 & \ldots & \ldots & \dots & \ldots \\
\hline
\mathbf 4 & 14 & \ldots & \ldots & \ldots & \dots & \ldots \\
\hline
\vdots & \vdots & \vdots & \vdots & \vdots&\ldots & \ddots\\
\end{array}
\]\noindent
Dit patroon heet \emph{de zigzagmethode van Cantor}. Daarmee wordt
$\Nat \times \Nat$ als volgt opgesomd:
\[
\tuple{0,0}, \tuple{0,1}, \tuple{1,0}, \tuple{0,2}, \tuple{1,1},
\tuple{2,0}, \tuple{0,3}, \tuple{1,2}, \tuple{2,1}, \tuple{3,0}, \dots
\]
Hieruit volgt:
\end{explain}

\begin{prop}\ollabel{natsquaredenumerable}
De verzameling $\Nat \times \Nat$ is !!{enumerable}.
\end{prop}

\begin{proof}
Zij $f \colon \Nat \to \Nat\times\Nat$ de functie die elke $k \in \Nat$
afbeeldt op het paar $\tuple{n,m} \in \Nat \times \Nat$ waarvoor $k$ de
waarde in rij $n$ en kolom $m$ van Cantors zigzagtabel is. 
\end{proof}

\begin{explain}
Deze techniek laat zich ook mooi uitbreiden. Zo kunnen we er de verzameling
geordende drietallen van natuurlijke getallen mee opsommen, namelijk:
\[
\Nat \times \Nat \times \Nat = \Setabs{\tuple{n,m,k}}{n,m,k \in \Nat}
\]
We beschouwen $\Nat \times \Nat \times \Nat$ als het cartesische product van
$\Nat \times \Nat$ met $\Nat$, dus
\[
\Nat^3 = (\Nat \times \Nat) \times \Nat =
\Setabs{\tuple{\tuple{n,m},k}}{n, m, k
  \in \Nat }
\]
Daarom kunnen we $\Nat^3$ met een tabel opsommen door op de ene as de
opsomming van $\Nat$ en op de andere as de opsomming van $\Nat^2$ te zetten:
\[
\begin{array}{ c | c | c | c | c | c}
& \mathbf 0 & \mathbf 1 & \mathbf 2 & \mathbf 3 & \dots \\
\hline
\mathbf{\tuple{0,0}} & \tuple{0,0,0} & \tuple{0,0,1} & \tuple{0,0,2} & \tuple{0,0,3} & \dots \\
\hline
\mathbf{\tuple{0,1}} & \tuple{0,1,0} & \tuple{0,1,1} & \tuple{0,1,2} & \tuple{0,1,3} & \dots \\
\hline
\mathbf{\tuple{1,0}} & \tuple{1,0,0} & \tuple{1,0,1} & \tuple{1,0,2} & \tuple{1,0,3} & \dots \\
\hline
\mathbf{\tuple{0,2}} & \tuple{0,2,0} & \tuple{0,2,1} & \tuple{0,2,2} & \tuple{0,2,3} & \dots\\
\hline
\vdots & \vdots & \vdots & \vdots & \vdots & \ddots \\
\end{array}
\]
Zo geeft een methode als Cantors zigzagmethode ook een opsomming van~$\Nat^3$.
Door zo verder te gaan, krijgen we opsommingen van $\Nat^n$ voor elk natuurlijk
getal $n$. Dus:
\end{explain}

\begin{prop}
$\Nat^n$ is !!{enumerable} voor elke $n \in \Nat$.
\end{prop}

\begin{prob}\label{sfr:siz:zigzag:prob:posint-n}
Bewijs dat $(\PosInt)^n$ !!{enumerable} is voor elke $n \in \Nat$.
\end{prob}

\begin{prob}\label{sfr:siz:zigzag:prob:posint-star}
  Bewijs dat $(\PosInt)^*$ !!{enumerable} is. Hierbij mag
  \cref{sfr:siz:zigzag:prob:posint-n} worden aangenomen.
\end{prob}
\end{document}